\documentclass[12pt,a4paper]{article}
\usepackage[top=2.5cm,bottom=2.5cm,left=2.5cm,right=2.5cm,headheight=15pt]{geometry}
\usepackage{amsmath,amssymb}
\usepackage{tikz}
\usetikzlibrary{shapes,arrows,arrows.meta,positioning,calc,backgrounds,fit,matrix}
\usepackage{booktabs,array,longtable,tabularx,multirow}
\usepackage{xcolor}
\usepackage{enumitem}
\usepackage{fancyhdr}
\usepackage{titlesec}
\usepackage{mdframed}
\usepackage{tcolorbox}
\tcbuselibrary{skins,breakable}
\usepackage{hyperref}

%% ---- Colors ----
\definecolor{folblue}{RGB}{0,82,165}
\definecolor{lightblue}{RGB}{220,235,255}
\definecolor{darkorange}{RGB}{180,70,0}
\definecolor{lightorange}{RGB}{255,235,210}
\definecolor{darkgreen}{RGB}{20,110,20}
\definecolor{lightgreen}{RGB}{200,240,205}
\definecolor{darkred}{RGB}{150,20,20}
\definecolor{lightred}{RGB}{255,210,210}
\definecolor{lightgray}{RGB}{242,242,242}

%% ---- tcolorbox styles ----
\newtcolorbox{qbox}[2][]{enhanced,breakable,
  colback=lightorange,colframe=darkorange,arc=4pt,
  title={\textbf{#2}},fonttitle=\bfseries\color{white},
  attach boxed title to top left={yshift=-2mm,xshift=4mm},
  boxed title style={colback=darkorange,arc=3pt},#1}

\newtcolorbox{solbox}[2][]{enhanced,breakable,
  colback=lightblue,colframe=folblue,arc=4pt,
  title={\textbf{#2}},fonttitle=\bfseries\color{white},
  attach boxed title to top left={yshift=-2mm,xshift=4mm},
  boxed title style={colback=folblue,arc=3pt},#1}

\newtcolorbox{keybox}[2][]{enhanced,breakable,
  colback=lightgreen,colframe=darkgreen,arc=4pt,
  title={\textbf{#2}},fonttitle=\bfseries\color{white},
  attach boxed title to top left={yshift=-2mm,xshift=4mm},
  boxed title style={colback=darkgreen,arc=3pt},#1}

\newtcolorbox{algobox}[2][]{enhanced,breakable,
  colback=lightgray,colframe=gray!60,arc=4pt,
  title={\textbf{#2}},fonttitle=\bfseries,
  attach boxed title to top left={yshift=-2mm,xshift=4mm},
  boxed title style={colback=gray!60,arc=3pt},#1}

\newtcolorbox{warnbox}[2][]{enhanced,breakable,
  colback=lightred,colframe=darkred,arc=4pt,
  title={\textbf{#2}},fonttitle=\bfseries\color{white},
  attach boxed title to top left={yshift=-2mm,xshift=4mm},
  boxed title style={colback=darkred,arc=3pt},#1}

%% ---- Page style ----
\pagestyle{fancy}\fancyhf{}
\lhead{\color{folblue}\textbf{RSA --- CSE 717 InfoSec}}
\rhead{\color{folblue}Exam: 17 Jun 2026}
\cfoot{--- \thepage\ ---}
\renewcommand{\headrulewidth}{0.4pt}

%% ---- Section style ----
\titleformat{\section}{\large\bfseries\color{folblue}}{}{0em}{}[\titlerule]
\titleformat{\subsection}{\normalsize\bfseries\color{darkorange}}{}{0em}{}

\begin{document}

\begin{center}
{\LARGE\bfseries\color{folblue} CSE 717 --- Information Security}\\[4pt]
{\large\bfseries Topic 8: RSA --- Key Generation \& Encryption/Decryption Numericals}\\[2pt]
{\large Complete Past Paper Solutions: 2020--2024}\\[6pt]
{\normalsize Source: Lc\#6 pp16--21 (RSA Algorithm Steps + worked examples)}\\[2pt]
{\normalsize\color{darkred}\textbf{HIGHEST-YIELD TOPIC: 5/5 years, 5--16 marks. Every Topic-8 question 2020--2024 included. Zero omissions.}}
\end{center}

\tableofcontents

%% ==========================================
\section*{Quick Reference --- The RSA Pipeline (CHEAT SHEET)}
%% ==========================================

\begin{warnbox}{CONFIRMED EVERY YEAR --- builds directly on Topic 7 (Extended Euclidean + Euler's totient)}
Every past-paper RSA question is the SAME 8-step pipeline with different small numbers. ``Given $e,n$ find $d$'' (2020 \& 2021, $e=31,n=3599$) is a near-certain repeat style. Always use \textbf{square-and-multiply} for $M^e \bmod n$ --- never multiply out the full power.
\end{warnbox}

\begin{algobox}{The 8-Step RSA Algorithm (Lc\#6 p16)}
\begin{enumerate}[noitemsep]
  \item Select two primes $p \neq q$.
  \item Compute $n = p \times q$.
  \item Compute $\phi(n) = (p-1)(q-1)$.
  \item Select $e$ such that $\gcd(e,\phi(n))=1$ and $1 < e < \phi(n)$.
  \item Compute $d = e^{-1} \bmod \phi(n)$, i.e.\ $ed \equiv 1 \pmod{\phi(n)}$ (Extended Euclidean --- Topic 7 Cheat Sheet B).
  \item \textbf{Public key} $= \{e,n\}$ (shared with everyone). \textbf{Private key} $= \{d,n\}$ (kept secret).
  \item \textbf{Encrypt:} $C = M^e \bmod n$, where $M < n$.
  \item \textbf{Decrypt:} $M = C^d \bmod n$.
\end{enumerate}

\textbf{Worked example from slides (Lc\#6 pp17--19): $p=13,q=11$}
\[
n=143,\quad \phi(n)=12\times10=120,\quad e=13\ (\gcd(13,120)=1)
\]
\[
13d \equiv 1 \pmod{120} \Rightarrow d=37 \quad(\text{check: } 13\times37=481=4\times120+1)
\]
\[
\text{Public key}=\{13,143\},\quad\text{Private key}=\{37,143\}
\]
\[
\text{Encrypt } P=13:\quad C=13^{13}\bmod143=52
\]
\[
\text{Decrypt } C=52:\quad P=52^{37}\bmod143=13 \checkmark
\]
\end{algobox}

\begin{keybox}{Square-and-Multiply --- the ONLY way to compute $M^e \bmod n$ by hand}
Write $e$ in binary. Scan bits left to right, starting with result $=1$:
\begin{itemize}[noitemsep]
  \item Always: \textbf{square} the result mod $n$.
  \item If the current bit is 1: also \textbf{multiply} by $M$ mod $n$.
\end{itemize}
Example: $e=17=10001_2$ needs only 4 squarings + 1 multiply (instead of 16 multiplications).
\end{keybox}

%% ==========================================
\section{2024 --- Q2(c): Encrypt ``DATA'' (3); Q2(d): Decrypt $c=17$ (3)}
%% ==========================================

\begin{qbox}{Question --- 2024 Q2(c) \& Q2(d)}
\begin{enumerate}[label=(\alph*)]
  \item[(c)] Encrypt the message ``DATA'' using the RSA cryptosystem with $p=61$, $q=53$, and $e=17$. \hfill [3]
  \item[(d)] An RSA-encrypted message consists of the ciphertext block $c=17$. The RSA system parameters are $p=61$, $q=53$, $e=17$. Decrypt the ciphertext and find the plaintext letter. \hfill [3]
\end{enumerate}
\end{qbox}

\subsection{Part (c): Encrypt ``DATA''}
\begin{solbox}{Answer --- 2024 Q2(c)}
\textbf{Step 1--3:} $n = p\times q = 61\times53 = \boxed{3233}$, \quad $\phi(n)=(60)(52)=\boxed{3120}$.

\textbf{Step 4:} $e=17$, check $\gcd(17,3120)=1$ \checkmark

\textbf{Encoding:} use each letter's ASCII code as the plaintext block $M$ (each $M<3233$, so each letter is its own block). $D=68,\,A=65,\,T=84,\,A=65$.

\textbf{Compute $C=M^{17}\bmod 3233$ via square-and-multiply} ($17=10001_2$ $\to$ 4 squarings, multiply by $M$ on bits 1 and 5):

\renewcommand{\arraystretch}{1.2}
\begin{center}
\begin{tabular}{lccccc}
\toprule
Letter ($M$) & after bit 1 & after bit 2 (sq) & after bit 3 (sq) & after bit 4 (sq) & after bit 5 (sq+$\times M$) \\
\midrule
D (68) & 68 & 1391 & 1547 & 789 & $\boxed{1759}$ \\
A (65) & 65 & 992 & 1232 & 1547 & $\boxed{2790}$ \\
T (84) & 84 & 590 & 2169 & 546 & $\boxed{2159}$ \\
A (65) & 65 & 992 & 1232 & 1547 & $\boxed{2790}$ \\
\bottomrule
\end{tabular}
\end{center}

\textbf{Ciphertext blocks:} $\text{``DATA''} \to \boxed{1759,\ 2790,\ 2159,\ 2790}$.

\textit{(This is the standard Stallings textbook example: $A=65 \to C=2790$.)}
\end{solbox}

\subsection{Part (d): Decrypt $c=17$}
\begin{solbox}{Answer --- 2024 Q2(d)}
\textbf{Find $d$:} need $d = e^{-1}\bmod\phi(n) = 17^{-1}\bmod 3120$ (Extended Euclidean, Topic 7 Cheat Sheet B).
\[
17\times2753 = 46801 = 15\times3120 + 1 \;\Rightarrow\; d=\boxed{2753}
\]

\textbf{Decrypt:} $M = C^d \bmod n = 17^{2753}\bmod3233$. Using repeated square-and-multiply on the 12-bit exponent $2753=101011000001_2$:
\[
M = 17^{2753}\bmod 3233 = \boxed{3170}
\]

\textbf{Note on ``plaintext letter'':} $3170$ falls outside the printable ASCII range ($0$--$255$), so this particular block does not map back to a single readable character under the ASCII encoding used in part (c). \textbf{The mark is for the correct method} (compute $d$ via Extended Euclidean, then $C^d\bmod n$ via square-and-multiply) --- show all steps even if the final number isn't a "nice" letter. If the examiner expects a letter, state the numeric result $3170$ and note it lies outside the A--Z/ASCII range.
\end{solbox}

%% ==========================================
\section{2023 --- Section B Q7(c): Encrypt ``CU'' (3); Q7(d): Decrypt (2)}
%% ==========================================

\begin{qbox}{Question --- 2023 Section B Q7(c) \& Q7(d)}
\begin{enumerate}[label=(\alph*)]
  \item[(c)] Encrypt the message ``CU'' using the RSA cryptosystem with $p=43$ and $q=59$. \hfill [3]
  \item[(d)] Consider an encrypted message ``0981 0461''. What is the decrypted message if it was encrypted using the RSA cryptosystem with $p=43$ and $q=59$? \hfill [2]
\end{enumerate}
\end{qbox}

\begin{warnbox}{Note: $e$ is not legible in the scanned paper for this question}
The value of $e$ is missing from the source PDF for Q7(c)/(d) (likely a scan/crop issue). $n$ and $\phi(n)$ depend only on $p,q$ and are fixed; $e$ must still satisfy $\gcd(e,\phi(n))=1$. Below, the \textbf{full pipeline} is demonstrated with a self-chosen valid $e=13$ so the worked example is internally consistent and exam-ready --- the \emph{method} is identical regardless of which valid $e$ the original paper used.
\end{warnbox}

\subsection{Part (c): Encrypt ``CU''}
\begin{solbox}{Answer --- 2023 Section B Q7(c)}
\textbf{Steps 1--3:} $n=p\times q=43\times59=\boxed{2537}$, \quad $\phi(n)=(42)(58)=\boxed{2436}$.

\textbf{Step 4:} choose $e=13$. Check $\gcd(13,2436)=1$ \checkmark (a valid public exponent).

\textbf{Encoding:} ASCII codes $C=67,\,U=85$ (each $<2537$, one block per letter).

\textbf{Compute $C=M^{13}\bmod 2537$ via square-and-multiply} ($13=1101_2$ $\to$ bits: 1,1,0,1):

\renewcommand{\arraystretch}{1.2}
\begin{center}
\begin{tabular}{lcccc}
\toprule
Letter ($M$) & after bit 1 & after bit 2 (sq+$\times M$) & after bit 3 (sq) & after bit 4 (sq+$\times M$) \\
\midrule
C (67) & 67 & 1397 & 656 & $\boxed{2044}$ \\
U (85) & 85 & 171 & 1334 & $\boxed{1246}$ \\
\bottomrule
\end{tabular}
\end{center}

\textbf{Ciphertext:} $\text{``CU''} \to \boxed{2044,\ 1246}$ \hfill (with $e=13,\ n=2537$)
\end{solbox}

\subsection{Part (d): Decrypt the ciphertext}
\begin{solbox}{Answer --- 2023 Section B Q7(d)}
\textbf{Find $d$:} $d=e^{-1}\bmod\phi(n)=13^{-1}\bmod2436$ (Extended Euclidean):
\[
13\times937 = 12181 = 5\times2436+1 \;\Rightarrow\; d=\boxed{937}
\]

\textbf{Decrypt} $M=C^d\bmod n=C^{937}\bmod2537$ for each block:
\[
2044^{937}\bmod2537 = \boxed{67} = \text{`C'}, \qquad 1246^{937}\bmod2537 = \boxed{85} = \text{`U'}
\]
$\Rightarrow$ plaintext recovered $= \text{``CU''}$ \checkmark

\textit{The paper's actual ciphertext ``0981 0461'' was produced with the original (unknown) $e$ -- the same Extended-Euclidean-then-square-and-multiply procedure applies verbatim once $e$ (and hence $d$) is known. Demonstrating the full round-trip (encrypt then decrypt back to the same letters) earns full method marks.}
\end{solbox}

%% ==========================================
\section{2022 --- Q6: Encrypt $M=3$, choose $e<10$ (3+3+2+8=16)}
%% ==========================================

\begin{qbox}{Question --- 2022 Q6}
Encrypt the message $M=3$ via RSA with $p=13$, $q=7$. Choose an $e<10$ and show all steps. Explain why an eavesdropper (Charlie), who intercepts the ciphertext, cannot decrypt the message without the private key. \hfill [3+3+2+8=16]
\end{qbox}

\begin{solbox}{Answer --- 2022 Q6}
\textbf{Steps 1--3:} $n=p\times q=13\times7=\boxed{91}$, \quad $\phi(n)=(12)(6)=\boxed{72}$.

\textbf{Step 4 --- choose $e<10$:} need $\gcd(e,72)=1$ with $1<e<10$. Candidates: $e=5$ ($\gcd(5,72)=1$ \checkmark) or $e=7$ ($\gcd(7,72)=1$ \checkmark). Take $\boxed{e=5}$.

\textbf{Step 5 --- find $d$:} $d=5^{-1}\bmod72$.
\[
5\times29=145=2\times72+1 \;\Rightarrow\; d=\boxed{29}
\]

\textbf{Step 6:} Public key $=\{5,91\}$, Private key $=\{29,91\}$.

\textbf{Step 7 --- encrypt $M=3$:} $C=3^5\bmod91$. $5=101_2$:
\[
3 \xrightarrow{\text{sq}} 9 \xrightarrow{\text{sq}\times3} 61 \quad\Rightarrow\quad C=\boxed{61}
\]

\textbf{Step 8 --- decrypt (verification):} $M=61^{29}\bmod91=\boxed{3}$ \checkmark

\vspace{4pt}
\textbf{Why Charlie can't decrypt without the private key:}
\begin{itemize}[noitemsep]
  \item Charlie sees the public key $\{e,n\}=\{5,91\}$ and the ciphertext $C=61$, but \textbf{not} $d=29$.
  \item Decryption requires $M=C^d\bmod n$ --- without $d$, Charlie must either:
  \begin{itemize}[noitemsep]
    \item \textbf{Factor $n=91$} into $p=13,q=7$ to recompute $\phi(n)$ and then $d$ --- for small textbook numbers this is feasible, but for real RSA ($n$ with 300+ digit primes) factoring is \textbf{computationally infeasible} (this is the RSA hardness assumption, based on the integer factorization problem).
    \item \textbf{Brute-force every possible $M<n$}, computing $M^e\bmod n$ and comparing to $C$ --- exponential search space for real key sizes.
  \end{itemize}
  \item Even knowing $e$ and $n$, computing $d=e^{-1}\bmod\phi(n)$ requires knowing $\phi(n)$, which requires knowing $p$ and $q$ --- and $n=pq$ does not reveal $p,q$ without factoring.
  \item \textbf{Conclusion:} the security of RSA rests entirely on the fact that factoring a large $n$ into its two prime factors is practically impossible with current computing power, so $d$ remains secret even though $e$ and $n$ are public.
\end{itemize}
\end{solbox}

%% ==========================================
\section{2021 \& 2020 --- $p=3,q=11,e=7,M=2$ (Encrypt+Decrypt); $e=31,n=3599$ find $d$}
%% ==========================================

\begin{qbox}{Question --- 2021 Section B Q2(a)\&(b), and identical 2020 Q6(a)\&(b)}
\begin{enumerate}[label=(\alph*)]
  \item Encrypt and decrypt using RSA with $p=3,q=11,e=7,M=2$. \hfill [2020: 4.75]\ [2021: 2.75]
  \item Given $e=31,\ n=3599$, find the private key $d$. \hfill [2+2 (both years)]
\end{enumerate}
\end{qbox}

\subsection{Part (a): RSA with $p=3,q=11,e=7,M=2$}
\begin{solbox}{Answer --- 2021 B2(a) / 2020 Q6(a)}
\textbf{Steps 1--3:} $n=3\times11=\boxed{33}$, \quad $\phi(n)=(2)(10)=\boxed{20}$.

\textbf{Step 4:} $e=7$, check $\gcd(7,20)=1$ \checkmark

\textbf{Step 5 --- find $d$:} $d=7^{-1}\bmod20$.
\[
7\times3=21=1\times20+1 \;\Rightarrow\; d=\boxed{3}
\]

\textbf{Step 6:} Public key $=\{7,33\}$, Private key $=\{3,33\}$.

\textbf{Step 7 --- encrypt $M=2$:} $C=2^7\bmod33$. $7=111_2$:
\[
2 \xrightarrow{\text{sq}\times2} 8 \xrightarrow{\text{sq}\times2} 29 \quad\Rightarrow\quad C=\boxed{29}
\]

\textbf{Step 8 --- decrypt:} $M=29^3\bmod33=\boxed{2}$ \checkmark
\end{solbox}

\subsection{Part (b): Given $e=31,\ n=3599$, find $d$}
\begin{solbox}{Answer --- 2021 B2(b) / 2020 Q6(b)}
\textbf{Factor $n=3599$} into its two primes (try small primes upward):
\[
3599 = 59 \times 61 \quad\Rightarrow\quad p=59,\ q=61
\]

\textbf{Compute $\phi(n)$:}
\[
\phi(3599)=(59-1)(61-1)=58\times60=\boxed{3480}
\]

\textbf{Find $d=31^{-1}\bmod3480$} via Extended Euclidean:
\[
31\times3031 = 93961 = 27\times3480+1 \;\Rightarrow\; d=\boxed{3031}
\]

\textbf{Private key} $=\{3031,\,3599\}$. \quad\textit{(This exact pair $e=31,n=3599$ repeats verbatim in 2020 \& 2021 --- memorize $n=59\times61$, $\phi(n)=3480$, $d=3031$.)}
\end{solbox}

%% ==========================================
\section*{Exam Strategy Summary}
%% ==========================================

\begin{keybox}{High-Yield Checklist for RSA (5/5 years, 5--16 marks --- LARGEST single-topic block some years)}
\begin{enumerate}[noitemsep]
  \item \textbf{Always small primes} in the range $\{3,7,11,13,17,31,43,53,59,61\}$ --- never need a calculator for $n=p\times q$ or $\phi(n)=(p-1)(q-1)$.
  \item \textbf{The 8-step pipeline never changes:} $n \to \phi(n) \to$ pick/verify $e \to d=e^{-1}\bmod\phi(n)$ (Extended Euclidean, Topic 7) $\to C=M^e\bmod n \to M=C^d\bmod n$ (square-and-multiply).
  \item \textbf{``Given $e,n$ find $d$''} is a near-certain repeat: memorize $e=31,n=3599 \Rightarrow p=59,q=61,\phi(n)=3480,d=3031$ (2020 \& 2021 verbatim).
  \item \textbf{Encryption of a word/letters} (2023 ``CU'', 2024 ``DATA''): encode each letter as its ASCII code, encrypt block-by-block with $C=M^e\bmod n$ --- one square-and-multiply per letter.
  \item \textbf{``Choose $e$''} questions (2022): pick the SMALLEST valid $e$ with $\gcd(e,\phi(n))=1$ --- usually $e=5$ or $e=7$ for small $\phi(n)$.
  \item \textbf{Conceptual add-on} (2022, 8 marks): "why can't an eavesdropper decrypt without the private key" $\to$ answer = \textbf{factoring $n$ is computationally infeasible} for large primes (RSA hardness assumption).
  \item \textbf{Square-and-multiply is mandatory} for any exponent $>$ a few --- write $e$ (or $d$) in binary, square-then-conditionally-multiply bit by bit. Never attempt to compute $M^e$ directly.
\end{enumerate}
\textit{This topic is built directly on Topic 7 (Extended Euclidean for $d$, Euler's totient for $\phi(n)$). If Topic 7 is solid, RSA is just "plug numbers into the same formula 8 times."}
\end{keybox}

\end{document}
